%=============================================================================
% Chapter 1: Rings and Ideals
%=============================================================================
\chapter{Rings and Ideals}\label{chapter1}

We shall begin by reviewing rapidly the definition and elementary properties of rings. This will indicate how much we are going to assume of the reader and it will also serve to fix notation and conventions. After this review we pass on to a discussion of prime and maximal ideals. The remainder of the chapter is devoted to explaining the various elementary operations which can be performed on ideals. The Grothendieck language of schemes is dealt with in the exercises at the end.


%=============================================================================
\section*{Rings and Ring Homomorphisms}
\addcontentsline{toc}{section}{Rings and Ring Homomorphisms}

A \emph{ring}\index{ring} $A$ is a set with two binary operations---(addition and multiplication) such that
\begin{enumerate}
    \item $A$ is an abelian group with respect to addition (so that $A$ has a zero element, denoted by $0$, and every $x \in A$ has an (additive) inverse, $-x$).
    \item Multiplication is associative $((xy)z = x(yz))$ and distributive over addition $(x(y + z) = xy + xz$, $(y + z)x = yx + zx)$.
\end{enumerate}

We shall consider only rings which are \emph{commutative}:
\begin{enumerate}
    \setcounter{enumi}{2}
    \item $xy = yx$ for all $x, y \in A$,
\end{enumerate}
and have an \emph{identity element} (denoted by $1$):
\begin{enumerate}
    \setcounter{enumi}{3}
    \item $\exists 1 \in A$ such that $x1 = 1x = x$ for all $x \in A$.
\end{enumerate}

The identity element is then unique.

Throughout this book the word ``ring'' shall mean a \textbf{commutative ring with an identity element}, that is, a ring satisfying axioms (1) to (4) above.

\begin{remark}
We do not exclude the possibility in (4) that $1$ might be equal to $0$. If so, then for any $x \in A$ we have
\[
    x = x1 = x0 = 0
\]
and so $A$ has only one element, $0$. In this case $A$ is the \emph{zero ring}\index{zero ring}, denoted by $0$ (by abuse of notation).
\end{remark}

A \emph{ring homomorphism}\index{ring homomorphism} is a mapping $f$ of a ring $A$ into a ring $B$ such that
\begin{enumerate}[i)]
    \item $f(x + y) = f(x) + f(y)$ (so that $f$ is a homomorphism of abelian groups, and therefore also $f(x - y) = f(x) - f(y)$, $f(-x) = -f(x)$, $f(0) = 0$),
    \item $f(xy) = f(x)f(y)$,
    \item $f(1) = 1$.
\end{enumerate}

In other words, $f$ respects addition, multiplication and the identity element.

A subset $S$ of a ring $A$ is a \emph{subring}\index{subring} of $A$ if $S$ is closed under addition and multiplication and contains the identity element of $A$. The identity mapping of $S$ into $A$ is then a ring homomorphism.

If $f \colon A \to B$, $g \colon B \to C$ are ring homomorphisms then so is their composition $g \circ f \colon A \to C$.


%=============================================================================
\section*{Ideals. Quotient Rings}
\addcontentsline{toc}{section}{Ideals. Quotient Rings}

An \emph{ideal}\index{ideal} $\mathfrak{a}$ of a ring $A$ is a subset of $A$ which is an additive subgroup and is such that $A\mathfrak{a} \subseteq \mathfrak{a}$ (i.e., $x \in A$ and $y \in \mathfrak{a}$ imply $xy \in \mathfrak{a}$). The quotient group $A/\mathfrak{a}$ inherits a uniquely defined multiplication from $A$ which makes it into a ring, called the \emph{quotient ring}\index{quotient ring} (or \emph{residue-class ring}\index{residue-class ring}) $A/\mathfrak{a}$. The elements of $A/\mathfrak{a}$ are the cosets of $\mathfrak{a}$ in $A$, and the mapping $\varphi \colon A \to A/\mathfrak{a}$ which maps each $x \in A$ to its coset $x + \mathfrak{a}$ is a surjective ring homomorphism.

We shall frequently use the following fact:

\begin{proposition}\label{prop:ideal-correspondence}
There is a one-to-one order-preserving correspondence between the ideals $\mathfrak{b}$ of $A$ which contain $\mathfrak{a}$, and the ideals $\bar{\mathfrak{b}}$ of $A/\mathfrak{a}$, given by $\mathfrak{b} = \varphi^{-1}(\bar{\mathfrak{b}})$. \qed
\end{proposition}

If $f \colon A \to B$ is any ring homomorphism, the \emph{kernel}\index{kernel} of $f$ ($= f^{-1}(0)$) is an ideal $\mathfrak{a}$ of $A$, and the \emph{image}\index{image} of $f$ ($= f(A)$) is a subring $C$ of $B$; and $f$ induces a ring isomorphism $A/\mathfrak{a} \cong C$.

We shall sometimes use the notation $x \equiv y \pmod{\mathfrak{a}}$; this means that $x - y \in \mathfrak{a}$.


%=============================================================================
\section*{Zero-Divisors. Nilpotent Elements. Units}
\addcontentsline{toc}{section}{Zero-Divisors. Nilpotent Elements. Units}

A \emph{zero-divisor}\index{zero-divisor} in a ring $A$ is an element $x$ which ``divides $0$'', i.e., for which there exists $y \neq 0$ in $A$ such that $xy = 0$. A ring with no zero-divisors $\neq 0$ (and in which $1 \neq 0$) is called an \emph{integral domain}\index{integral domain}. For example, $\Z$ and $k[x_1, \ldots, x_n]$ ($k$ a field, $x_i$ indeterminates) are integral domains.

An element $x \in A$ is \emph{nilpotent}\index{nilpotent} if $x^n = 0$ for some $n > 0$. A nilpotent element is a zero-divisor (unless $A = 0$), but not conversely (in general).

A \emph{unit}\index{unit} in $A$ is an element $x$ which ``divides $1$'', i.e., an element $x$ such that $xy = 1$ for some $y \in A$. The element $y$ is then uniquely determined by $x$, and is written $x^{-1}$. The units in $A$ form a (multiplicative) abelian group.

The multiples $ax$ of an element $x \in A$ form a \emph{principal ideal}\index{principal ideal}, denoted by $(x)$ or $Ax$. $x$ is a unit $\iff (x) = A = (1)$. The zero ideal $(0)$ is usually denoted by $0$.

A \emph{field}\index{field} is a ring $A$ in which $1 \neq 0$ and every non-zero element is a unit. Every field is an integral domain (but not conversely: $\Z$ is not a field).

\begin{proposition}\label{prop:field-characterization}
Let $A$ be a ring $\neq 0$. Then the following are equivalent:
\begin{enumerate}[i)]
    \item $A$ is a field;
    \item the only ideals in $A$ are $0$ and $(1)$;
    \item every homomorphism of $A$ into a non-zero ring $B$ is injective.
\end{enumerate}
\end{proposition}

\begin{proof}
\noindent i) $\implies$ ii). Let $\mathfrak{a} \neq 0$ be an ideal in $A$. Then $\mathfrak{a}$ contains a non-zero element $x$; $x$ is a unit, hence $\mathfrak{a} \supseteq (x) = (1)$, hence $\mathfrak{a} = (1)$.

ii) $\implies$ iii). Let $\phi \colon A \to B$ be a ring homomorphism. Then $\Ker(\phi)$ is an ideal $\neq (1)$ in $A$, hence $\Ker(\phi) = 0$, hence $\phi$ is injective.

iii) $\implies$ i). Let $x$ be an element of $A$ which is not a unit. Then $(x) \neq (1)$, hence $B = A/(x)$ is not the zero ring. Let $\phi \colon A \to B$ be the natural homomorphism of $A$ onto $B$, with kernel $(x)$. By hypothesis, $\phi$ is injective, hence $(x) = 0$, hence $x = 0$.
\end{proof}


%=============================================================================
\section*{Prime Ideals and Maximal Ideals}
\addcontentsline{toc}{section}{Prime Ideals and Maximal Ideals}

An ideal $\mathfrak{p}$ in $A$ is \emph{prime}\index{prime ideal} if $\mathfrak{p} \neq (1)$ and if $xy \in \mathfrak{p} \implies x \in \mathfrak{p}$ or $y \in \mathfrak{p}$.

An ideal $\mathfrak{m}$ in $A$ is \emph{maximal}\index{maximal ideal} if $\mathfrak{m} \neq (1)$ and if there is no ideal $\mathfrak{a}$ such that $\mathfrak{m} \subset \mathfrak{a} \subset (1)$ (strict inclusions). Equivalently:
\[
    \mathfrak{p} \text{ is prime } \iff A/\mathfrak{p} \text{ is an integral domain};
\]
\[
    \mathfrak{m} \text{ is maximal } \iff A/\mathfrak{m} \text{ is a field (by \eqref{prop:ideal-correspondence} and \eqref{prop:field-characterization})}.
\]
Hence a maximal ideal is prime (but not conversely, in general). The zero ideal is prime $\iff A$ is an integral domain.

If $f \colon A \to B$ is a ring homomorphism and $\mathfrak{q}$ is a prime ideal in $B$, then $f^{-1}(\mathfrak{q})$ is a prime ideal in $A$, for $A/f^{-1}(\mathfrak{q})$ is isomorphic to a subring of $B/\mathfrak{q}$ and hence has no zero-divisor $\neq 0$. But if $\mathfrak{n}$ is a maximal ideal of $B$ it is not necessarily true that $f^{-1}(\mathfrak{n})$ is maximal in $A$; all we can say for sure is that it is prime. (Example: $A = \Z$, $B = \Q$, $\mathfrak{n} = 0$.)

Prime ideals are fundamental to the whole of commutative algebra. The following theorem and its corollaries ensure that there is always a sufficient supply of them.

\begin{theorem}\label{thm:maximal-ideal-exists}
Every ring $A \neq 0$ has at least one maximal ideal. (Remember that ``ring'' means commutative ring with $1$.)
\end{theorem}

\begin{proof}
This is a standard application of Zorn's lemma.\footnote{Let $S$ be a non-empty partially ordered set (i.e., we are given a relation $x \leq y$ on $S$ which is reflexive and transitive and such that $x \leq y$ and $y \leq x$ together imply $x = y$). A subset $T$ of $S$ is a \emph{chain} if either $x \leq y$ or $y \leq x$ for every pair of elements $x, y$ in $T$. Then Zorn's lemma may be stated as follows: if every chain $T$ of $S$ has an upper bound in $S$ (i.e., if there exists $x \in S$ such that $t \leq x$ for all $t \in T$) then $S$ has at least one maximal element.

For a proof of the equivalence of Zorn's lemma with the axiom of choice, the well-ordering principle, etc., see for example P. R. Halmos, \emph{Naive Set Theory}, Van Nostrand (1960).} Let $\Sigma$ be the set of all ideals $\neq (1)$ in $A$. Order $\Sigma$ by inclusion. $\Sigma$ is not empty, since $0 \in \Sigma$. To apply Zorn's lemma we must show that every chain in $\Sigma$ has an upper bound in $\Sigma$; let then $(\mathfrak{a}_\alpha)$ be a chain of ideals in $\Sigma$, so that for each pair of indices $\alpha, \beta$ we have either $\mathfrak{a}_\alpha \subseteq \mathfrak{a}_\beta$ or $\mathfrak{a}_\beta \subseteq \mathfrak{a}_\alpha$. Let $\mathfrak{a} = \bigcup_\alpha \mathfrak{a}_\alpha$. Then $\mathfrak{a}$ is an ideal (verify this) and $1 \notin \mathfrak{a}$ because $1 \notin \mathfrak{a}_\alpha$ for all $\alpha$. Hence $\mathfrak{a} \in \Sigma$, and $\mathfrak{a}$ is an upper bound of the chain. Hence by Zorn's lemma $\Sigma$ has a maximal element.
\end{proof}

\begin{corollary}\label{cor:ideal-in-maximal}
If $\mathfrak{a} \neq (1)$ is an ideal of $A$, there exists a maximal ideal of $A$ containing $\mathfrak{a}$.
\end{corollary}

\begin{proof}
Apply \eqref{thm:maximal-ideal-exists} to $A/\mathfrak{a}$, bearing in mind \eqref{prop:ideal-correspondence}. Alternatively, modify the proof of \eqref{thm:maximal-ideal-exists}.
\end{proof}

\begin{corollary}\label{cor:nonunit-in-maximal}
Every non-unit of $A$ is contained in a maximal ideal. \qed
\end{corollary}

\begin{remark}
1) If $A$ is Noetherian (Chapter 7) we can avoid the use of Zorn's lemma: the set of all ideals $\neq (1)$ has a maximal element.

2) There exist rings with exactly one maximal ideal, for example fields. A ring $A$ with exactly one maximal ideal $\mathfrak{m}$ is called a \emph{local ring}\index{local ring}. The field $k = A/\mathfrak{m}$ is called the \emph{residue field}\index{residue field} of $A$.
\end{remark}

\begin{proposition}\label{prop:local-ring-characterization}
\begin{enumerate}[i)]
    \item Let $A$ be a ring and $\mathfrak{m} \neq (1)$ an ideal of $A$ such that every $x \in A \setminus \mathfrak{m}$ is a unit in $A$. Then $A$ is a local ring and $\mathfrak{m}$ its maximal ideal.
    \item Let $A$ be a ring and $\mathfrak{m}$ a maximal ideal of $A$, such that every element of $1 + \mathfrak{m}$ (i.e., every $1 + x$, where $x \in \mathfrak{m}$) is a unit in $A$. Then $A$ is a local ring.
\end{enumerate}
\end{proposition}

\begin{proof}
\noindent i) Every ideal $\neq (1)$ consists of non-units, hence is contained in $\mathfrak{m}$. Hence $\mathfrak{m}$ is the only maximal ideal of $A$.

ii) Let $x \in A \setminus \mathfrak{m}$. Since $\mathfrak{m}$ is maximal, the ideal generated by $x$ and $\mathfrak{m}$ is $(1)$, hence there exist $y \in A$ and $t \in \mathfrak{m}$ such that $xy + t = 1$; hence $xy = 1 - t$ belongs to $1 + \mathfrak{m}$ and therefore is a unit. Now use i).
\end{proof}

A ring with only a finite number of maximal ideals is called \emph{semi-local}\index{semi-local ring}.

\begin{examples}
    \begin{enumerate}[1)]
        \item $A = k[x_1, \ldots, x_n]$, $k$ a field. Let $f \in A$ be an irreducible polynomial. By unique factorization, the ideal $(f)$ is prime.
        \item $A = \Z$. Every ideal in $\Z$ is of the form $(m)$ for some $m \geq 0$. The ideal $(m)$ is prime $\iff m = 0$ or a prime number. All the ideals $(p)$, where $p$ is a prime number, are maximal: $\Z/(p)$ is the field of $p$ elements.
        
        The same holds in Example 1) for $n = 1$, but not for $n > 1$. The ideal $\mathfrak{m}$ of all polynomials in $A = k[x_1, \ldots, x_n]$ with zero constant term is maximal (since it is the kernel of the homomorphism $A \to k$ which maps $f \in A$ to $f(0)$). But if $n > 1$, $\mathfrak{m}$ is not a principal ideal: in fact it requires at least $n$ generators.
        \item \label{ex:principal-ideal-domain}A \emph{principal ideal domain}\index{principal ideal domain} is an integral domain in which every ideal is principal. In such a ring every non-zero prime ideal is maximal. For if $(x) \neq 0$ is a prime ideal and $(y) \supset (x)$, we have $x \in (y)$, say $x = yz$, so that $yz \in (x)$ and $y \notin (x)$, hence $z \in (x)$: say $z = tx$. Then $x = yz = ytx$, so that $yt = 1$ and therefore $(y) = (1)$.
    \end{enumerate}
\end{examples}


%=============================================================================
\section*{Nilradical and Jacobson Radical}
\addcontentsline{toc}{section}{Nilradical and Jacobson Radical}

\begin{proposition}\label{prop:nilradical-ideal}
The set $\nilrad$ of all nilpotent elements in a ring $A$ is an ideal, and $A/\nilrad$ has no nilpotent element $\neq 0$.
\end{proposition}

\begin{proof}
If $x \in \nilrad$, clearly $ax \in \nilrad$ for all $a \in A$. Let $x, y \in \nilrad$: say $x^m = 0$, $y^n = 0$. By the binomial theorem (which is valid in any commutative ring), $(x + y)^{m+n-1}$ is a sum of integer multiples of products $x^r y^s$, where $r + s = m + n - 1$; we cannot have both $r < m$ and $s < n$, hence each of these products vanishes and therefore $(x + y)^{m+n-1} = 0$. Hence $x + y \in \nilrad$ and therefore $\nilrad$ is an ideal.

Let $\bar{x} \in A/\nilrad$ be represented by $x \in A$. Then $\bar{x}^n$ is represented by $x^n$, so that $\bar{x}^n = 0 \implies x^n \in \nilrad \implies (x^n)^k = 0$ for some $k > 0 \implies x \in \nilrad \implies \bar{x} = 0$.
\end{proof}

The ideal $\nilrad$ is called the \emph{nilradical}\index{nilradical} of $A$. The following proposition gives an alternative definition of $\nilrad$:

\begin{proposition}\label{prop:nilradical-intersection}
The nilradical of $A$ is the intersection of all the prime ideals of $A$.
\end{proposition}

\begin{proof}
Let $\nilrad'$ denote the intersection of all the prime ideals of $A$. If $f \in A$ is nilpotent and if $\mathfrak{p}$ is a prime ideal, then $f^n = 0 \in \mathfrak{p}$ for some $n > 0$, hence $f \in \mathfrak{p}$ (because $\mathfrak{p}$ is prime). Hence $f \in \nilrad'$.

Conversely, suppose that $f$ is not nilpotent. Let $\Sigma$ be the set of ideals $\mathfrak{a}$ with the property
\[
    n > 0 \implies f^n \notin \mathfrak{a}.
\]
Then $\Sigma$ is not empty because $0 \in \Sigma$. As in \eqref{thm:maximal-ideal-exists} Zorn's lemma can be applied to the set $\Sigma$, ordered by inclusion, and therefore $\Sigma$ has a maximal element. Let $\mathfrak{p}$ be a maximal element of $\Sigma$. We shall show that $\mathfrak{p}$ is a prime ideal. Let $x, y \notin \mathfrak{p}$. Then the ideals $\mathfrak{p} + (x)$, $\mathfrak{p} + (y)$ strictly contain $\mathfrak{p}$ and therefore do not belong to $\Sigma$; hence
\[
    f^m \in \mathfrak{p} + (x), \quad f^n \in \mathfrak{p} + (y)
\]
for some $m, n$. It follows that $f^{m+n} \in \mathfrak{p} + (xy)$, hence the ideal $\mathfrak{p} + (xy)$ is not in $\Sigma$ and therefore $xy \notin \mathfrak{p}$. Hence we have a prime ideal $\mathfrak{p}$ such that $f \notin \mathfrak{p}$, so that $f \notin \nilrad'$.
\end{proof}

The \emph{Jacobson radical}\index{Jacobson radical} $\jacrad$ of $A$ is defined to be the intersection of all the maximal ideals of $A$. It can be characterized as follows:

\begin{proposition}\label{prop:jacobson-radical}
$x \in \jacrad \iff 1 - xy$ is a unit in $A$ for all $y \in A$.
\end{proposition}

\begin{proof}
\noindent $\implies$: Suppose $1 - xy$ is not a unit. By \eqref{cor:nonunit-in-maximal} it belongs to some maximal ideal $\mathfrak{m}$; but $x \in \jacrad \subseteq \mathfrak{m}$, hence $xy \in \mathfrak{m}$ and therefore $1 \in \mathfrak{m}$, which is absurd.

$\impliedby$: Suppose $x \notin \mathfrak{m}$ for some maximal ideal $\mathfrak{m}$. Then $\mathfrak{m}$ and $x$ generate the unit ideal $(1)$, so that we have $u + xy = 1$ for some $u \in \mathfrak{m}$ and some $y \in A$. Hence $1 - xy \in \mathfrak{m}$ and is therefore not a unit.
\end{proof}


%=============================================================================
\section*{Operations on Ideals}
\addcontentsline{toc}{section}{Operations on Ideals}

If $\mathfrak{a}$, $\mathfrak{b}$ are ideals in a ring $A$, their \emph{sum}\index{ideal!sum} $\mathfrak{a} + \mathfrak{b}$ is the set of all $x + y$ where $x \in \mathfrak{a}$ and $y \in \mathfrak{b}$. It is the smallest ideal containing $\mathfrak{a}$ and $\mathfrak{b}$. More generally, we may define the sum $\sum_{i \in I} \mathfrak{a}_i$ of any family (possibly infinite) of ideals $\mathfrak{a}_i$ of $A$; its elements are all sums $\sum x_i$, where $x_i \in \mathfrak{a}_i$ for all $i \in I$ and almost all of the $x_i$ (i.e., all but a finite set) are zero. It is the smallest ideal of $A$ which contains all the ideals $\mathfrak{a}_i$.

The \emph{intersection}\index{ideal!intersection} of any family $(\mathfrak{a}_i)_{i \in I}$ of ideals is an ideal. Thus the ideals of $A$ form a complete lattice with respect to inclusion.

The \emph{product}\index{ideal!product} of two ideals $\mathfrak{a}$, $\mathfrak{b}$ in $A$ is the ideal $\mathfrak{a}\mathfrak{b}$ generated by all products $xy$, where $x \in \mathfrak{a}$ and $y \in \mathfrak{b}$. It is the set of all finite sums $\sum x_i y_i$ where each $x_i \in \mathfrak{a}$ and each $y_i \in \mathfrak{b}$. Similarly we define the product of any finite family of ideals. In particular the powers $\mathfrak{a}^n$ ($n > 0$) of an ideal $\mathfrak{a}$ are defined; conventionally, $\mathfrak{a}^0 = (1)$. Thus $\mathfrak{a}^n$ ($n > 0$) is the ideal generated by all products $x_1 x_2 \cdots x_n$ in which each factor $x_i$ belongs to $\mathfrak{a}$.

\begin{examples}
1) If $A = \Z$, $\mathfrak{a} = (m)$, $\mathfrak{b} = (n)$ then $\mathfrak{a} + \mathfrak{b}$ is the ideal generated by the $\hcf$ of $m$ and $n$; $\mathfrak{a} \cap \mathfrak{b}$ is the ideal generated by their $\lcm$; and $\mathfrak{a}\mathfrak{b} = (mn)$. Thus (in this case) $\mathfrak{a}\mathfrak{b} = \mathfrak{a} \cap \mathfrak{b} \iff m, n$ are coprime.

2) $A = k[x_1, \ldots, x_n]$, $\mathfrak{a} = (x_1, \ldots, x_n) =$ ideal generated by $x_1, \ldots, x_n$. Then $\mathfrak{a}^m$ is the set of all polynomials with no terms of degree $< m$.
\end{examples}

The three operations so far defined (sum, intersection, product) are all commutative and associative. Also there is the distributive law
\[
    \mathfrak{a}(\mathfrak{b} + \mathfrak{c}) = \mathfrak{a}\mathfrak{b} + \mathfrak{a}\mathfrak{c}.
\]

In the ring $\Z$, $\cap$ and $+$ are distributive over each other. This is not the case in general, and the best we have in this direction is the \emph{modular law}\index{modular law}
\[
    \mathfrak{a} \cap (\mathfrak{b} + \mathfrak{c}) = \mathfrak{a} \cap \mathfrak{b} + \mathfrak{a} \cap \mathfrak{c} \quad \text{if } \mathfrak{a} \supseteq \mathfrak{b} \text{ or } \mathfrak{a} \supseteq \mathfrak{c}.
\]

Again, in $\Z$, we have $(\mathfrak{a} + \mathfrak{b})(\mathfrak{a} \cap \mathfrak{b}) = \mathfrak{a}\mathfrak{b}$; but in general we have only $(\mathfrak{a} + \mathfrak{b})(\mathfrak{a} \cap \mathfrak{b}) \subseteq \mathfrak{a}\mathfrak{b}$ (since $(\mathfrak{a} + \mathfrak{b})(\mathfrak{a} \cap \mathfrak{b}) = \mathfrak{a}(\mathfrak{a} \cap \mathfrak{b}) + \mathfrak{b}(\mathfrak{a} \cap \mathfrak{b}) \subseteq \mathfrak{a}\mathfrak{b}$). Clearly $\mathfrak{a}\mathfrak{b} \subseteq \mathfrak{a} \cap \mathfrak{b}$, hence
\[
    \mathfrak{a} \cap \mathfrak{b} = \mathfrak{a}\mathfrak{b} \quad \text{provided } \mathfrak{a} + \mathfrak{b} = (1).
\]

Two ideals $\mathfrak{a}, \mathfrak{b}$ are said to be \emph{coprime} (or \emph{comaximal}) if $\mathfrak{a} + \mathfrak{b} = (1)$. Thus for coprime ideals we have $\mathfrak{a} \cap \mathfrak{b} = \mathfrak{a}\mathfrak{b}$. Clearly two ideals $\mathfrak{a}, \mathfrak{b}$ are coprime if and only if there exist $x \in \mathfrak{a}$ and $y \in \mathfrak{b}$ such that $x + y = 1$.

Let $A_1, \ldots, A_n$ be rings. Their \emph{direct product}
\[
A = \prod_{i=1}^{n} A_i
\]
is the set of all sequences $x = (x_1, \ldots, x_n)$ with $x_i \in A_i$ $(1 \le i \le n)$ and componentwise addition and multiplication. $A$ is a commutative ring with identity element $(1, 1, \ldots, 1)$. We have projections $p_i \colon A \to A_i$ defined by $p_i(x) = x_i$; they are ring homomorphisms.

Let $A$ be a ring and $\mathfrak{a}_1, \ldots, \mathfrak{a}_n$ ideals of $A$. Define a homomorphism
\[
\phi \colon A \to \prod_{i=1}^{n} (A/\mathfrak{a}_i)
\]
by the rule $\phi(x) = (x + \mathfrak{a}_1, \ldots, x + \mathfrak{a}_n)$.

\begin{proposition}\label{pro:chinese-remainder}
\begin{enumerate}
    \item[i)] If $\mathfrak{a}_i, \mathfrak{a}_j$ are coprime whenever $i \neq j$, then $\prod \mathfrak{a}_i = \bigcap \mathfrak{a}_i$.
    \item[ii)] $\phi$ is surjective $\iff$ $\mathfrak{a}_i, \mathfrak{a}_j$ are coprime whenever $i \neq j$.
    \item[iii)] $\phi$ is injective $\iff \bigcap \mathfrak{a}_i = (0)$.
\end{enumerate}
\end{proposition}

\begin{proof}
i) by induction on $n$. The case $n = 2$ is dealt with above. Suppose $n > 2$ and the result true for $\mathfrak{a}_1, \ldots, \mathfrak{a}_{n-1}$, and let $\mathfrak{b} = \prod_{i=1}^{n-1} \mathfrak{a}_i = \bigcap_{i=1}^{n-1} \mathfrak{a}_i$. Since $\mathfrak{a}_i + \mathfrak{a}_n = (1)$ $(1 \le i \le n-1)$ we have equations $x_i + y_i = 1$ $(x_i \in \mathfrak{a}_i, y_i \in \mathfrak{a}_n)$ and therefore
\[
\prod_{i=1}^{n-1} x_i = \prod_{i=1}^{n-1} (1 - y_i) \equiv 1 \pmod{\mathfrak{a}_n}.
\]
Hence $\mathfrak{a}_n + \mathfrak{b} = (1)$ and so
\[
\prod_{i=1}^{n} \mathfrak{a}_i = \mathfrak{b}\mathfrak{a}_n = \mathfrak{b} \cap \mathfrak{a}_n = \bigcap_{i=1}^{n} \mathfrak{a}_i.
\]

ii) $\implies$: Let us show for example that $\mathfrak{a}_1, \mathfrak{a}_2$ are coprime. There exists $x \in A$ such that $\phi(x) = (1, 0, \ldots, 0)$; hence $x \equiv 1 \pmod{\mathfrak{a}_1}$ and $x \equiv 0 \pmod{\mathfrak{a}_2}$, so that
\[
1 = (1 - x) + x \in \mathfrak{a}_1 + \mathfrak{a}_2.
\]

$\impliedby$: It is enough to show, for example, that there is an element $x \in A$ such that $\phi(x) = (1, 0, \ldots, 0)$. Since $\mathfrak{a}_1 + \mathfrak{a}_i = (1)$ $(i > 1)$ we have equations $u_i + v_i = 1$ $(u_i \in \mathfrak{a}_1, v_i \in \mathfrak{a}_i)$. Take $x = \prod_{i=2}^{n} v_i$, then $x = \prod (1 - u_i) \equiv 1 \pmod{\mathfrak{a}_1}$, and $x \equiv 0 \pmod{\mathfrak{a}_i}, i > 1$. Hence $\phi(x) = (1, 0, \ldots, 0)$ as required.

iii) Clear, since $\bigcap \mathfrak{a}_i$ is the kernel of $\phi$.
\end{proof}

The \emph{union} $\mathfrak{a} \cup \mathfrak{b}$ of ideals is not in general an ideal.

\begin{proposition}\label{pro:prime avoidance}
\begin{enumerate}
    \item[i)] Let $\mathfrak{p}_1, \ldots, \mathfrak{p}_n$ be prime ideals and let $\mathfrak{a}$ be an ideal contained in $\bigcup_{i=1}^{n} \mathfrak{p}_i$. Then $\mathfrak{a} \subseteq \mathfrak{p}_i$ for some $i$.
    \item[ii)] Let $\mathfrak{a}_1, \ldots, \mathfrak{a}_n$ be ideals and let $\mathfrak{p}$ be a prime ideal containing $\bigcap_{i=1}^{n} \mathfrak{a}_i$. Then $\mathfrak{p} \supseteq \mathfrak{a}_i$ for some $i$. If $\mathfrak{p} = \bigcap \mathfrak{a}_i$, then $\mathfrak{p} = \mathfrak{a}_i$ for some $i$.
\end{enumerate}
\end{proposition}

\begin{proof}
i) is proved by induction on $n$ in the form
\[
\mathfrak{a} \not\subseteq \mathfrak{p}_i\ (1 \le i \le n) \implies \mathfrak{a} \not\subseteq \bigcup_{i=1}^{n} \mathfrak{p}_i.
\]
It is certainly true for $n = 1$. If $n > 1$ and the result is true for $n - 1$, then for each $i$ there exists $x_i \in \mathfrak{a}$ such that $x_i \notin \mathfrak{p}_j$ whenever $j \neq i$. If for some $i$ we have $x_i \notin \mathfrak{p}_i$, we are through. If not, then $x_i \in \mathfrak{p}_i$ for all $i$. Consider the element
\[
y = \sum_{i=1}^{n} x_1 x_2 \cdots x_{i-1} x_{i+1} x_{i+2} \cdots x_n;
\]
we have $y \in \mathfrak{a}$ and $y \notin \mathfrak{p}_i\ (1 \le i \le n)$. Hence $\mathfrak{a} \not\subseteq \bigcup_{i=1}^{n} \mathfrak{p}_i$.

ii) Suppose $\mathfrak{p} \not\supseteq \mathfrak{a}_i$ for all $i$. Then there exist $x_i \in \mathfrak{a}_i, x_i \notin \mathfrak{p}\ (1 \le i \le n)$, and therefore $\prod x_i \in \prod \mathfrak{a}_i \subseteq \bigcap \mathfrak{a}_i$; but $\prod x_i \notin \mathfrak{p}$ (since $\mathfrak{p}$ is prime). Hence $\mathfrak{p} \not\supseteq \bigcap \mathfrak{a}_i$. Finally, if $\mathfrak{p} = \bigcap \mathfrak{a}_i$, then $\mathfrak{p} \subseteq \mathfrak{a}_i$, and hence $\mathfrak{p} = \mathfrak{a}_i$ for some $i$.
\end{proof}

If $\mathfrak{a}, \mathfrak{b}$ are ideals in a ring $A$, their \emph{ideal quotient} is
\[
(\mathfrak{a} : \mathfrak{b}) = \{x \in A : x\mathfrak{b} \subseteq \mathfrak{a}\}
\]
which is an ideal. In particular, $(0 : \mathfrak{b})$ is called the \emph{annihilator} of $\mathfrak{b}$ and is also denoted by $\operatorname{Ann}(\mathfrak{b})$: it is the set of all $x \in A$ such that $x\mathfrak{b} = 0$. In this notation the set of all zero-divisors in $A$ is
\[
D = \bigcup_{x \neq 0} \operatorname{Ann}(x).
\]
If $\mathfrak{b}$ is a principal ideal $(x)$, we shall write $(\mathfrak{a} : x)$ in place of $(\mathfrak{a} : (x))$.

\begin{example}
    If $A = \mathbb{Z}$, $\mathfrak{a} = (m)$, $\mathfrak{b} = (n)$, where say $m = \prod_p p^{\mu_p}$, $n = \prod_p p^{\nu_p}$, then $(\mathfrak{a} : \mathfrak{b}) = (q)$ where $q = \prod_p p^{\gamma_p}$ and
    \[
    \gamma_p = \max(\mu_p - \nu_p, 0) = \mu_p - \min(\mu_p, \nu_p).
    \]
    Hence $q = m/(m, n)$, where $(m, n)$ is the h.c.f.\ of $m$ and $n$.
\end{example}

\begin{exercise}
    \begin{enumerate}
    \item[i)] $\mathfrak{a} \subseteq (\mathfrak{a} : \mathfrak{b})$
    \item[ii)] $(\mathfrak{a} : \mathfrak{b})\mathfrak{b} \subseteq \mathfrak{a}$
    \item[iii)] $((\mathfrak{a} : \mathfrak{b}) : \mathfrak{c}) = (\mathfrak{a} : \mathfrak{b}\mathfrak{c}) = ((\mathfrak{a} : \mathfrak{c}) : \mathfrak{b})$
    \item[iv)] $\left(\bigcap_i \mathfrak{a}_i : \mathfrak{b}\right) = \bigcap_i (\mathfrak{a}_i : \mathfrak{b})$
    \item[v)] $(\mathfrak{a} : \sum_i \mathfrak{b}_i) = \bigcap_i (\mathfrak{a} : \mathfrak{b}_i)$.
\end{enumerate}
\end{exercise}

If $\mathfrak{a}$ is any ideal of $A$, the \emph{radical} of $\mathfrak{a}$ is
\[
r(\mathfrak{a}) = \{x \in A : x^n \in \mathfrak{a} \text{ for some } n > 0\}.
\]
If $f \colon A \to A/\mathfrak{a}$ is the standard homomorphism, then $r(\mathfrak{a}) = f^{-1}(\mathfrak{N}_{A/\mathfrak{a}})$ and hence $r(\mathfrak{a})$ is an ideal by \eqref{prop:nilradical-ideal}.

\begin{exercise}\label{exer:radical-properties}
    \begin{enumerate}
    \item[i)] $r(\mathfrak{a}) \supseteq \mathfrak{a}$
    \item[ii)] $r(r(\mathfrak{a})) = r(\mathfrak{a})$
    \item[iii)] $r(\mathfrak{a}\mathfrak{b}) = r(\mathfrak{a} \cap \mathfrak{b}) = r(\mathfrak{a}) \cap r(\mathfrak{b})$
    \item[iv)] $r(\mathfrak{a}) = (1) \iff \mathfrak{a} = (1)$
    \item[v)] $r(\mathfrak{a} + \mathfrak{b}) = r(r(\mathfrak{a}) + r(\mathfrak{b}))$
    \item[vi)] if $\mathfrak{p}$ is prime, $r(\mathfrak{p}^n) = \mathfrak{p}$ for all $n > 0$.
\end{enumerate}
\end{exercise}

\begin{proposition}\label{prop:radical-intersection-primes}
The radical of an ideal $\mathfrak{a}$ is the intersection of the prime ideals which contain $\mathfrak{a}$.
\end{proposition}

\begin{proof}
    Apply \eqref{prop:nilradical-intersection} to $A/\mathfrak{a}$.
\end{proof}

More generally, we may define the radical $r(E)$ of any subset $E$ of $A$ in the same way. It is not an ideal in general. We have $r(\bigcup_\alpha E_\alpha) = \bigcup r(E_\alpha)$, for any family of subsets $E_\alpha$ of $A$.

\begin{proposition}\label{prop:zerodivisors-union}
    $D\coloneqq$ the set of zero-divisors of $A$ is $\bigcup_{x \neq 0} r(\Ann(x))$.
\end{proposition}

\begin{proof}
$D = r(D) = r(\bigcup_{x \neq 0} \Ann(x)) = \bigcup_{x \neq 0} r(\Ann(x))$.
\end{proof}

\begin{example}
    If $A = \mathbb{Z}$, $\mathfrak{a} = (m)$, let $p_i$ $(1 \le i \le r)$ be the distinct prime divisors of $m$. Then $r(\mathfrak{a}) = (p_1 \cdots p_r) = \bigcap_{i=1}^{r} (p_i)$.
\end{example}

\begin{proposition}\label{prop:coprime-radicals}
    Let $\mathfrak{a}$, $\mathfrak{b}$ be ideals in a ring $A$ such that $r(\mathfrak{a})$, $r(\mathfrak{b})$ are coprime. Then $\mathfrak{a}$, $\mathfrak{b}$ are coprime.
\end{proposition}

\begin{proof}
$r(\mathfrak{a} + \mathfrak{b}) = r(r(\mathfrak{a}) + r(\mathfrak{b})) = r((1)) = (1)$, hence $\mathfrak{a} + \mathfrak{b} = (1)$.
\end{proof}

\section*{Extension and Contraction}
\addcontentsline{toc}{section}{Extension and Contraction}

Let $f \colon A \to B$ be a ring homomorphism. If $\mathfrak{a}$ is an ideal in $A$, the set $f(\mathfrak{a})$ is not necessarily an ideal in $B$ (e.g., let $f$ be the embedding of $\mathbb{Z}$ in $\mathbb{Q}$, the field of rationals, and take $\mathfrak{a}$ to be any non-zero ideal in $\mathbb{Z}$.) We define the \emph{extension} $\mathfrak{a}^e$ of $\mathfrak{a}$ to be the ideal $Bf(\mathfrak{a})$ generated by $f(\mathfrak{a})$ in $B$: explicitly, $\mathfrak{a}^e$ is the set of all sums $\sum y_i f(x_i)$ where $x_i \in \mathfrak{a}$, $y_i \in B$.

If $\mathfrak{b}$ is an ideal of $B$, then $f^{-1}(\mathfrak{b})$ is always an ideal of $A$, called the \emph{contraction} $\mathfrak{b}^c$ of $\mathfrak{b}$. If $\mathfrak{b}$ is prime, then $\mathfrak{b}^c$ is prime. If $\mathfrak{a}$ is prime, $\mathfrak{a}^e$ need not be prime (for example, $f \colon \mathbb{Z} \to \mathbb{Q}$, $\mathfrak{a} \neq 0$; then $\mathfrak{a}^e = \mathbb{Q}$, which is not a prime ideal).

We can factorize $f$ as follows:
\[
A \stackrel{p}{\longrightarrow} f(A) \stackrel{j}{\longrightarrow} B
\]
where $p$ is surjective and $j$ is injective. For $p$ the situation is very simple \eqref{prop:ideal-correspondence}: there is a one-to-one correspondence between ideals of $f(A)$ and ideals of $A$ which contain $\Ker(f)$, and prime ideals correspond to prime ideals. For $j$, on the other hand, the general situation is very complicated. The classical example is from algebraic number theory.

\begin{example}
    Consider $\mathbb{Z} \to \mathbb{Z}[\mathrm{i}]$, where $i = \sqrt{-1}$. A prime ideal $(p)$ of $\mathbb{Z}$ may or may not stay prime when extended to $\mathbb{Z}[\mathrm{i}]$. In fact $\mathbb{Z}[\mathrm{i}]$ is a principal ideal domain (because it has a Euclidean algorithm) and the situation is as follows:

    \begin{enumerate}
        \item[i)] $(2)^e = ((1+\mathrm{i})^2)$, the \emph{square} of a prime ideal in $\mathbb{Z}[\mathrm{i}]$;
        \item[ii)] If $p \equiv 1 \pmod 4$ then $(p)^e$ is the product of two distinct prime ideals (for example, $(5)^e = (2+\mathrm{i})(2-\mathrm{i})$);
        \item[iii)] If $p \equiv 3 \pmod 4$ then $(p)^e$ is prime in $\mathbb{Z}[\mathrm{i}]$.
    \end{enumerate}

    Of these, ii) is not a trivial result. It is effectively equivalent to a theorem of Fermat which says that a prime $p \equiv 1 \pmod 4$ can be expressed, essentially uniquely, as a sum of two integer squares (thus $5 = 2^2 + 1^2$, $97 = 9^2 + 4^2$, etc.).
\end{example}

In fact the behavior of prime ideals under extensions of this sort is one of the central problems of algebraic number theory.

\begin{proposition}\label{prop:extension-contraction-properties}
Let $f \colon A \to B$, $\mathfrak{a}$ and $\mathfrak{b}$ be as before. Then:
\begin{enumerate}[i)]
    \item $\mathfrak{a} \subseteq \mathfrak{a}^{ec}$, $\mathfrak{b} \supseteq \mathfrak{b}^{ce}$;
    \item $\mathfrak{b}^c = \mathfrak{b}^{cec}$, $\mathfrak{a}^e = \mathfrak{a}^{ece}$;
    \item If $C$ is the set of contracted ideals in $A$ and if $E$ is the set of extended ideals in $B$, then $C = \{\mathfrak{a} : \mathfrak{a}^{ec} = \mathfrak{a}\}$, $E = \{\mathfrak{b} : \mathfrak{b}^{ce} = \mathfrak{b}\}$, and $\mathfrak{a} \mapsto \mathfrak{a}^e$ is a bijective map of $C$ onto $E$, whose inverse is $\mathfrak{b} \mapsto \mathfrak{b}^c$.
\end{enumerate}
\end{proposition}

\begin{proof}
i) is trivial, and ii) follows from i).

iii) If $\mathfrak{a} \in C$, then $\mathfrak{a} = \mathfrak{b}^c = \mathfrak{b}^{cec} = \mathfrak{a}^{ec}$; conversely if $\mathfrak{a} = \mathfrak{a}^{ec}$ then $\mathfrak{a}$ is the contraction of $\mathfrak{a}^e$. Similarly for $E$.
\end{proof}

\begin{exercise}\label{ex:extension-contraction-operations}
If $\mathfrak{a}_1, \mathfrak{a}_2$ are ideals of $A$ and if $\mathfrak{b}_1, \mathfrak{b}_2$ are ideals of $B$, then
\begin{align*}
    (\mathfrak{a}_1 + \mathfrak{a}_2)^e &= \mathfrak{a}_1^e + \mathfrak{a}_2^e, & (\mathfrak{b}_1 + \mathfrak{b}_2)^c &\supseteq \mathfrak{b}_1^c + \mathfrak{b}_2^c, \\
    (\mathfrak{a}_1 \cap \mathfrak{a}_2)^e &\subseteq \mathfrak{a}_1^e \cap \mathfrak{a}_2^e, & (\mathfrak{b}_1 \cap \mathfrak{b}_2)^c &= \mathfrak{b}_1^c \cap \mathfrak{b}_2^c, \\
    (\mathfrak{a}_1 \mathfrak{a}_2)^e &= \mathfrak{a}_1^e \mathfrak{a}_2^e, & (\mathfrak{b}_1 \mathfrak{b}_2)^c &\supseteq \mathfrak{b}_1^c \mathfrak{b}_2^c, \\
    (\mathfrak{a}_1 : \mathfrak{a}_2)^e &\subseteq (\mathfrak{a}_1^e : \mathfrak{a}_2^e), & (\mathfrak{b}_1 : \mathfrak{b}_2)^c &\subseteq (\mathfrak{b}_1^c : \mathfrak{b}_2^c), \\
    r(\mathfrak{a})^e &\subseteq r(\mathfrak{a}^e), & r(\mathfrak{b})^c &= r(\mathfrak{b}^c).
\end{align*}
The set of ideals $E$ is closed under sum and product, and $C$ is closed under the other three operations.
\end{exercise}

%=============================================================================
% Chapter 1 Exercises - Rewritten to match original order with hints
%=============================================================================
\section*{Exercises}
\addcontentsline{toc}{section}{Exercises}

\begin{enumerate}
    \item \label{exer1-chapter1}Let $x$ be a nilpotent element of a ring $A$. Show that $1 + x$ is a unit of $A$. Deduce that the sum of a nilpotent element and a unit is a unit.

    \item \label{exer2-chapter1}Let $A$ be a ring and let $A[x]$ be the ring of polynomials in an indeterminate $x$, with coefficients in $A$. Let $f = a_0 + a_1 x + \cdots + a_n x^n \in A[x]$. Prove that:
    \begin{enumerate}[i)]
        \item $f$ is a unit in $A[x]$ $\iff$ $a_0$ is a unit in $A$ and $a_1, \ldots, a_n$ are nilpotent.\\
        \textit{Hint:} If $b_0 + b_1 x + \cdots + b_m x^m$ is the inverse of $f$, prove by induction on $r$ that $a_n^{r+1} b_{m-r} = 0$. Hence show that $a_n$ is nilpotent, and then use Exercise \ref{exer1-chapter1}.
        \item $f$ is nilpotent $\iff$ $a_0, a_1, \ldots, a_n$ are nilpotent.
        \item $f$ is a zero-divisor $\iff$ there exists $a \neq 0$ in $A$ such that $af = 0$.\\
        \textit{Hint:} Choose a polynomial $g = b_0 + b_1 x + \cdots + b_m x^m$ of least degree $m$ such that $fg = 0$. Then $a_n b_m = 0$, hence $a_n g = 0$ (because $a_n g$ annihilates $f$ and has degree $< m$). Now show by induction that $a_{n-r} g = 0$ ($0 \leq r \leq n$).
        \item $f$ is said to be \emph{primitive} if $(a_0, a_1, \ldots, a_n) = (1)$. Prove that if $f, g \in A[x]$, then $fg$ is primitive $\iff$ $f$ and $g$ are primitive.
    \end{enumerate}

    \item \label{exer3-chapter1}Generalize the results of Exercise \ref{exer2-chapter1} to a polynomial ring $A[x_1, \ldots, x_r]$ in several indeterminates.

    \item \label{exer4-chapter1}In the ring $A[x]$, the Jacobson radical is equal to the nilradical.

    \item \label{exer5-chapter1}Let $A$ be a ring and let $A[[x]]$ be the ring of formal power series $f = \sum_{n=0}^\infty a_n x^n$ with coefficients in $A$. Show that:
    \begin{enumerate}[i)]
        \item $f$ is a unit in $A[[x]]$ $\iff$ $a_0$ is a unit in $A$.
        \item If $f$ is nilpotent, then $a_n$ is nilpotent for all $n \geq 0$. Is the converse true? (See Chapter \ref{chapter7}, Exercise \ref{exer2-chapter7}.)
        \item $f$ belongs to the Jacobson radical of $A[[x]]$ $\iff$ $a_0$ belongs to the Jacobson radical of $A$.
        \item The contraction of a maximal ideal $\mathfrak{m}$ of $A[[x]]$ is a maximal ideal of $A$, and $\mathfrak{m}$ is generated by $\mathfrak{m}^c$ and $x$.
        \item Every prime ideal of $A$ is the contraction of a prime ideal of $A[[x]]$.
    \end{enumerate}

    \item \label{exer6-chapter1}A ring $A$ is such that every ideal not contained in the nilradical contains a non-zero idempotent (that is, an element $e$ such that $e^2 = e \neq 0$). Prove that the nilradical and Jacobson radical of $A$ are equal.

    \item \label{exer7-chapter1}Let $A$ be a ring in which every element $x$ satisfies $x^n = x$ for some $n > 1$ (depending on $x$). Show that every prime ideal in $A$ is maximal.

    \item \label{exer8-chapter1}Let $A$ be a ring $\neq 0$. Show that the set of prime ideals of $A$ has minimal elements with respect to inclusion.

    \item \label{exer9-chapter1}Let $\mathfrak{a}$ be an ideal $\neq (1)$ in a ring $A$. Show that $\mathfrak{a} = r(\mathfrak{a})$ $\Rightarrow$ $\mathfrak{a}$ is an intersection of prime ideals.

    \item \label{exer10-chapter1}Let $A$ be a ring, $\nilrad$ its nilradical. Show that the following are equivalent:
    \begin{enumerate}[i)]
        \item $A$ has exactly one prime ideal;
        \item every element of $A$ is either a unit or nilpotent;
        \item $A/\nilrad$ is a field.
    \end{enumerate}

    \item \label{exer11-chapter1}A ring $A$ is \emph{Boolean} if $x^2 = x$ for all $x \in A$. In a Boolean ring $A$, show that:
    \begin{enumerate}[i)]
        \item $2x = 0$ for all $x \in A$;
        \item every prime ideal $\mathfrak{p}$ is maximal, and $A/\mathfrak{p}$ is a field with two elements;
        \item every finitely generated ideal in $A$ is principal.
    \end{enumerate}

    \item \label{exer12-chapter1}A local ring contains no idempotent $\neq 0, 1$.

    \emph{Construction of an algebraic closure of a field}(E. Artin). 
    \item \label{exer13-chapter1}Let $K$ be a field and let $\Sigma$ be the set of all irreducible monic polynomials $f$ in one indeterminate with coefficients in $K$. Let $A$ be the polynomial ring over $K$ generated by indeterminates $x_f$, one for each $f \in \Sigma$. Let $\mathfrak{a}$ be the ideal of $A$ generated by the polynomials $f(x_f)$ for all $f \in \Sigma$. Show that $\mathfrak{a} \neq (1)$.
    
    \quad\, Let $\mathfrak{m}$ be a maximal ideal of $A$ containing $\mathfrak{a}$, and let $K_1 = A/\mathfrak{m}$. Then $K_1$ is an extension field of $K$ in which each $f \in \Sigma$ has a root. Repeat the construction with $K_1$ in place of $K$, obtaining a field $K_2$, and so on. Let $L = \bigcup_{n=1}^\infty K_n$. Then $L$ is a field in which each $f \in \Sigma$ splits completely into linear factors. Let $\bar{K}$ be the set of all elements of $L$ which are algebraic over $K$. Then $\bar{K}$ is an algebraic closure of $K$.

    \item \label{exer14-chapter1}In a ring $A$, let $\Sigma$ be the set of all ideals in which every element is a zero-divisor. Show that the set $\Sigma$ has maximal elements and that every maximal element of $\Sigma$ is a prime ideal. Hence the set of zero-divisors in $A$ is a union of prime ideals.

    \emph{The prime spectrum of a ring.}
    \item \label{exer15-chapter1}Let $A$ be a ring and let $X$ be the set of all prime ideals of $A$. For each subset $E$ of $A$, let $V(E)$ denote the set of all prime ideals of $A$ which contain $E$. Prove that:
    \begin{enumerate}[i)]
        \item if $\mathfrak{a}$ is the ideal generated by $E$, then $V(E) = V(\mathfrak{a}) = V(r(\mathfrak{a}))$.
        \item $V(0) = X$, $V(1) = \varnothing$.
        \item if $(E_i)_{i \in I}$ is any family of subsets of $A$, then $V(\bigcup_{i \in I} E_i) = \bigcap_{i \in I} V(E_i)$.
        \item $V(\mathfrak{a} \cap \mathfrak{b}) = V(\mathfrak{a}\mathfrak{b}) = V(\mathfrak{a}) \cup V(\mathfrak{b})$ for any ideals $\mathfrak{a}, \mathfrak{b}$ of $A$.
    \end{enumerate}
    These results show that the sets $V(E)$ satisfy the axioms for closed sets in a topological space. The resulting topology is called the \emph{Zariski topology}\index{Zariski, ring!topology}. The topological space $X$ is called the \emph{prime spectrum}\index{Spectrum, maximal!prime} of $A$, and is written $\Spec(A)$.

    \item \label{exer16-chapter1}Draw pictures of $\Spec(\Z)$, $\Spec(\R)$, $\Spec(\C[x])$, $\Spec(\R[x])$, $\Spec(\Z[x])$.

    \item \label{exer17-chapter1}For each $f \in A$, let $X_f$ denote the complement of $V(f)$ in $X = \Spec(A)$. The sets $X_f$ are open. Show that they form a basis of open sets for the Zariski topology, and that:
    \begin{enumerate}[i)]
        \item $X_f \cap X_g = X_{fg}$;
        \item $X_f = \varnothing$ $\iff$ $f$ is nilpotent;
        \item $X_f = X$ $\iff$ $f$ is a unit;
        \item $X_f = X_g$ $\iff$ $r((f)) = r((g))$;
        \item $X$ is quasi-compact (that is, every open covering of $X$ has a finite sub-covering).\\
        \textit{Hint:} To prove (v), remark that it is enough to consider a covering of $X$ by basic open sets $X_{f_i}$ ($i \in I$). Show that the $f_i$ generate the unit ideal and hence that there is an equation of the form $1 = \sum_{i \in J} g_i f_i$ ($g_i \in A$) where $J$ is some finite subset of $I$. Then the $X_{f_i}$ ($i \in J$) cover $X$.
        \item More generally, each $X_f$ is quasi-compact.
        \item An open subset of $X$ is quasi-compact if and only if it is a finite union of sets $X_f$.
    \end{enumerate}
    The sets $X_f$ are called \emph{basic open sets} of $X = \Spec(A)$.

    \item \label{exer18-chapter1}For psychological reasons it is sometimes convenient to denote a prime ideal of $A$ by a letter such as $x$ or $y$ when thinking of it as a point of $X = \Spec(A)$. When thinking of $x$ as a prime ideal of $A$, we denote it by $\mathfrak{p}_x$ (logically, of course, it is the same thing). Show that:
    \begin{enumerate}[i)]
        \item the set $\{x\}$ is closed (we say that $x$ is a ``closed point'') in $\Spec(A)$ $\iff$ $\mathfrak{p}_x$ is maximal;
        \item $\overline{\{x\}} = V(\mathfrak{p}_x)$;
        \item $y \in \overline{\{x\}}$ $\iff$ $\mathfrak{p}_x \subseteq \mathfrak{p}_y$;
        \item $X$ is a $T_0$-space (this means that if $x, y$ are distinct points of $X$, then either there is a neighborhood of $x$ which does not contain $y$, or else there is a neighborhood of $y$ which does not contain $x$).
    \end{enumerate}

    \item \label{exer19-chapter1}A topological space $X$ is said to be \emph{irreducible} if $X \neq \varnothing$ and if every pair of non-empty open sets in $X$ intersect, or equivalently if every non-empty open set is dense in $X$. Show that $\Spec(A)$ is irreducible if and only if the nilradical of $A$ is a prime ideal.

    \item \label{exer20-chapter1}Let $X$ be a topological space.
    \begin{enumerate}[i)]
        \item If $Y$ is an irreducible subspace of $X$, then the closure $\bar{Y}$ of $Y$ in $X$ is irreducible.
        \item Every irreducible subspace of $X$ is contained in a maximal irreducible subspace.
        \item The maximal irreducible subspaces of $X$ are closed and cover $X$. They are called the \emph{irreducible components} of $X$. What are the irreducible components of a Hausdorff space?
        \item If $A$ is a ring and $X = \Spec(A)$, then the irreducible components of $X$ are the closed sets $V(\mathfrak{p})$, where $\mathfrak{p}$ is a minimal prime ideal of $A$ (Exercise \ref{exer8-chapter1}).
    \end{enumerate}

    \item \label{exer21-chapter1}Let $\varphi \colon A \to B$ be a ring homomorphism. Let $X = \Spec(A)$ and $Y = \Spec(B)$. If $\mathfrak{q} \in Y$, then $\varphi^{-1}(\mathfrak{q})$ is a prime ideal of $A$, i.e., a point of $X$. Hence $\varphi$ induces a mapping $\varphi^* \colon Y \to X$. Show that:
    \begin{enumerate}[i)]
        \item If $f \in A$ then $\varphi^{*-1}(X_f) = Y_{\varphi(f)}$, and hence that $\varphi^*$ is continuous.
        \item If $\mathfrak{a}$ is an ideal of $A$, then $\varphi^{*-1}(V(\mathfrak{a})) = V(\mathfrak{a}^e)$.
        \item If $\mathfrak{b}$ is an ideal of $B$, then $\varphi^*(V(\mathfrak{b})) = V(\mathfrak{b}^c)$.
        \item If $\varphi$ is surjective, then $\varphi^*$ is a homeomorphism of $Y$ onto the closed subset $V(\Ker(\varphi))$ of $X$. (In particular, $\Spec(A)$ and $\Spec(A/\nilrad)$ (where $\nilrad$ is the nilradical of $A$) are naturally homeomorphic.)
        \item If $\varphi$ is injective, then $\varphi^*(Y)$ is dense in $X$. More precisely, $\varphi^*(Y)$ is dense in $X$ $\iff$ $\Ker(\varphi) \subseteq \nilrad$.
        \item Let $\psi \colon B \to C$ be another ring homomorphism. Then $(\psi\varphi)^* = \varphi^* \circ \psi^*$.
        \item Let $A$ be an integral domain with just one non-zero prime ideal $\mathfrak{p}$, and let $K$ be the field of fractions of $A$. Let $B = (A/\mathfrak{p}) \times K$. Define $\varphi \colon A \to B$ by $\varphi(x) = (\bar{x}, x)$, where $\bar{x}$ is the image of $x$ in $A/\mathfrak{p}$. Show that $\varphi^*$ is bijective but not a homeomorphism.
    \end{enumerate}

    \item \label{exer22-chapter1}Let $A = \prod_{i=1}^n A_i$ be the direct product of rings $A_i$. Show that $\Spec(A)$ is the disjoint union of open (and closed) subspaces $X_i$, where $X_i$ is canonically homeomorphic with $\Spec(A_i)$.
    
    \quad\, Conversely, let $A$ be any ring. Show that the following statements are equivalent:
    \begin{enumerate}[i)]
        \item $X = \Spec(A)$ is disconnected.
        \item $A \cong A_1 \times A_2$ where neither of the rings $A_1, A_2$ is the zero ring.
        \item $A$ contains an idempotent $\neq 0, 1$.
    \end{enumerate}
    In particular, the spectrum of a local ring is always connected (Exercise \ref{exer12-chapter1}).

    \item \label{exer23-chapter1}Let $A$ be a Boolean ring\index{Boolean ring}\index{Ring!Boolean} (Exercise \ref{exer11-chapter1}), and let $X = \Spec(A)$.
    \begin{enumerate}[i)]
        \item For each $f \in A$, the set $X_f$ (Exercise \ref{exer17-chapter1}) is both open and closed in $X$.
        \item Let $f_1, \ldots, f_n \in A$. Show that $X_{f_1} \cup \cdots \cup X_{f_n} = X_f$ for some $f \in A$.
        \item The sets $X_f$ are the only subsets of $X$ which are both open and closed.
        
        \textit{Hint:} Let $Y \subseteq X$ be both open and closed. Since $Y$ is open, it is a union of basic open sets $X_f$. Since $Y$ is closed and $X$ is quasi-compact (Exercise \ref{exer17-chapter1}), $Y$ is quasi-compact. Hence $Y$ is a finite union of basic open sets; now use (ii) above.
        \item $X$ is a compact Hausdorff space.
    \end{enumerate}

    \item \label{exer24-chapter1}Let $L$ be a lattice, in which the sup and inf of two elements $a, b$ are denoted by $a \vee b$ and $a \wedge b$ respectively. $L$ is a \emph{Boolean lattice} (or \emph{Boolean algebra}) if:
    \begin{enumerate}[i)]
        \item $L$ has a least element and a greatest element (denoted by $0, 1$ respectively).
        \item Each of $\vee, \wedge$ is distributive over the other.
        \item Each $a \in L$ has a unique ``complement'' $a' \in L$ such that $a \vee a' = 1$ and $a \wedge a' = 0$.
    \end{enumerate}
    (For example, the set of all subsets of a set, ordered by inclusion, is a Boolean lattice.)

    \quad\, Let $L$ be a Boolean lattice. Define addition and multiplication in $L$ by the rules:
    \[
        a + b = (a \wedge b') \vee (a' \wedge b), \qquad ab = a \wedge b.
    \]
    Verify that in this way $L$ becomes a Boolean ring, say $A(L)$.
    
    \quad\, Conversely, starting from a Boolean ring $A$, define an ordering on $A$ as follows: $a \leq b$ means that $a = ab$. Show that, with respect to this ordering, $A$ is a Boolean lattice. [The sup and inf are given by $a \vee b = a + b + ab$ and $a \wedge b = ab$, and the complement by $a' = 1 - a$.] In this way we obtain a one-to-one correspondence between (isomorphism classes of) Boolean rings and (isomorphism classes of) Boolean lattices.

    \item \label{exer25-chapter1}From the last two exercises deduce \emph{Stone's theorem}, that every Boolean lattice is isomorphic to the lattice of open-and-closed subsets of some compact Hausdorff topological space.

    \item \label{exer26-chapter1}Let $A$ be a ring. The subspace of $\Spec(A)$ consisting of the maximal ideals of $A$, with the induced topology, is called the \emph{maximal spectrum}\index{Spectrum, maximal} of $A$ and is denoted by $\Max(A)$. For arbitrary commutative rings it does not have the nice functorial properties of $\Spec(A)$ (see Exercise \ref{exer21-chapter1}), because the inverse image of a maximal ideal under a ring homomorphism need not be maximal.
    
    \quad\, Let $X$ be a compact Hausdorff space and let $C(X)$ denote the ring of all real-valued continuous functions on $X$ (add and multiply functions by adding and multiplying their values). For each $x \in X$, let $\mathfrak{m}_x$ be the set of all $f \in C(X)$ such that $f(x) = 0$. The ideal $\mathfrak{m}_x$ is maximal, because it is the kernel of the (surjective) homomorphism $C(X) \to \R$ which takes $f$ to $f(x)$. If $\mathfrak{X}$ denotes $\Max(C(X))$, we have therefore defined a mapping $\mu \colon X \to \mathfrak{X}$, namely $x \mapsto \mathfrak{m}_x$.

    \quad\, We shall show that $\mu$ is a homeomorphism of $X$ onto $\mathfrak{X}$.
    \begin{enumerate}[i)]
        \item Let $\mathfrak{m}$ be any maximal ideal of $C(X)$, and let $V = V(\mathfrak{m})$ be the set of common zeros of the functions in $\mathfrak{m}$: that is, $V = \{x \in X : f(x) = 0 \text{ for all } f \in \mathfrak{m}\}$. Suppose that $V$ is empty. Then for each $x \in X$ there exists $f_x \in \mathfrak{m}$ such that $f_x(x) \neq 0$. Since $f_x$ is continuous, there is an open neighborhood $U_x$ of $x$ in $X$ on which $f_x$ does not vanish. By compactness a finite number of the neighborhoods, say $U_{x_1}, \ldots, U_{x_n}$, cover $X$. Let $f = f_{x_1}^2 + \cdots + f_{x_n}^2$. Then $f$ does not vanish at any point of $X$, hence is a unit in $C(X)$. But this contradicts $f \in \mathfrak{m}$, hence $V$ is not empty. Let $x$ be a point of $V$. Then $\mathfrak{m} \subseteq \mathfrak{m}_x$, hence $\mathfrak{m} = \mathfrak{m}_x$ because $\mathfrak{m}$ is maximal. Hence $\mu$ is surjective.
        \item By Urysohn's lemma (this is the only non-trivial fact required in the argument) the continuous functions separate the points of $X$. Hence $x \neq y$ $\Rightarrow$ $\mathfrak{m}_x \neq \mathfrak{m}_y$, and therefore $\mu$ is injective.
        \item Let $f \in C(X)$; let $U_f = \{x \in X : f(x) \neq 0\}$ and let $\mathcal{U}_f = \{\mathfrak{m} \in \mathfrak{X} : f \notin \mathfrak{m}\}$. Show that $\mu(U_f) = \mathcal{U}_f$. The open sets $U_f$ (resp. $\mathcal{U}_f$) form a basis of the topology of $X$ (resp. $\mathfrak{X}$) and therefore $\mu$ is a homeomorphism.
        
        Thus $X$ can be reconstructed from the ring of functions $C(X)$.
    \end{enumerate}
    

    \emph{Affine algebraic varieties.}\index{Varieties, affine algebraic}
    \item \label{exer27-chapter1} Let $k$ be an algebraically closed field and let $\{f_\alpha(t_1, \ldots, t_n) = 0\}$ be a set of polynomial equations in $n$ variables with coefficients in $k$. The set $X$ of all points $x = (x_1, \ldots, x_n) \in k^n$ which satisfy these equations is an affine algebraic variety.

    Consider the set of all polynomials $g \in k[t_1, \ldots, t_n]$ with the property that $g(x) = 0$ for all $x \in X$. This set is an ideal $I(X)$ in the polynomial ring, and is called the \emph{ideal of the variety} $X$. The quotient ring $P(X) = k[t_1, \ldots, t_n]/I(X)$ is the ring of polynomial functions on $X$, because two polynomials $g, h$ define the same polynomial function on $X$ if and only if $g - h$ vanishes at every point of $X$, that is, if and only if $g - h \in I(X)$.
    
    \quad\, Let $\xi_i$ be the image of $t_i$ in $P(X)$. The $\xi_i$ ($1 \leq i \leq n$) are the coordinate functions on $X$: if $x \in X$, then $\xi_i(x)$ is the $i$th coordinate of $x$. $P(X)$ is generated as a $k$-algebra by the coordinate functions, and is called the \emph{coordinate ring} (or \emph{affine algebra}) of $X$.

    \quad\, As in Exercise \ref{exer26-chapter1}, for each $x \in X$ let $\mathfrak{m}_x$ be the ideal of all $f \in P(X)$ such that $f(x) = 0$; it is a maximal ideal of $P(X)$. Hence, if $\tilde{X} = \Max(P(X))$, we have defined a mapping $\mu \colon X \to \tilde{X}$, namely $x \mapsto \mathfrak{m}_x$.

    It is easy to show that $\mu$ is injective: if $x \neq y$, we must have $x_i \neq y_i$ for some $i$ ($1 \leq i \leq n$), and hence $\xi_i - x_i$ is in $\mathfrak{m}_x$ but not in $\mathfrak{m}_y$, so that $\mathfrak{m}_x \neq \mathfrak{m}_y$. What is less obvious (but still true) is that $\mu$ is surjective. This is one form of \emph{Hilbert's Nullstellensatz} (see Chapter \ref{chapter7}).

    \item \label{exer28-chapter1}Let $f_1, \ldots, f_m$ be elements of $k[t_1, \ldots, t_n]$. They determine a polynomial mapping $\varphi \colon k^n \to k^m$: if $x \in k^n$, the coordinates of $\varphi(x)$ are $f_1(x), \ldots, f_m(x)$.

    \quad\, Let $X, Y$ be affine algebraic varieties in $k^n, k^m$ respectively. A mapping $\varphi \colon X \to Y$ is said to be \emph{regular} if $\varphi$ is the restriction to $X$ of a polynomial mapping from $k^n$ to $k^m$.

    \quad\, If $\eta$ is a polynomial function on $Y$, then $\eta \circ \varphi$ is a polynomial function on $X$. Hence $\varphi$ induces a $k$-algebra homomorphism $P(Y) \to P(X)$, namely $\eta \mapsto \eta \circ \varphi$. Show that in this way we obtain a one-to-one correspondence between the regular mappings $X \to Y$ and the $k$-algebra homomorphisms $P(Y) \to P(X)$.
\end{enumerate}

